\section{Introdução}

A entrega frequente de mudanças precisa ser analisada em conjunto com a
estabilidade do software. O trabalho de Forsgren, Humble e Kim
\cite{accelerate} e as métricas de desempenho de entrega do DORA \cite{dora}
motivam o estudo dessa relação. Neste laboratório, o interesse recai sobre
frequência de entrega, tempo entre a escrita e a entrega de mudanças,
proporção de entregas com falha e tempo de recuperação.

Repositórios públicos oferecem releases, commits e registros de integração
contínua, mas esses registros não demonstram, por si só, quando uma mudança
chegou à produção ou provocou um incidente. Uma release de biblioteca pode
ser disponibilizada sem adoção imediata; uma falha de CI pode impedir uma
entrega em vez de representar falha em produção. Os valores minerados são,
portanto, \textit{proxies}, cujo significado depende de suas definições
operacionais.

O objetivo deste estudo é caracterizar proxies DORA em projetos open-source
que utilizam GitHub Actions e avaliar a confiança nas comparações resultantes.
As questões de pesquisa investigam frequência de releases (RQ01), lead time
(RQ02), taxas de falha (RQ03), recuperação de workflows (RQ04), associação
entre frequência e falha (RQ05), diferenças entre características dos projetos
(RQ06) e sensibilidade da classificação às definições operacionais (RQ07).

O protocolo da disciplina prevê uma janela de doze meses, uma amostra final
de pelo menos 300 repositórios elegíveis e definições comuns. Para lead time,
serão comparadas a mediana dos tempos por release, considerando o commit
mais antigo de cada entrega, e a mediana dos tempos de todos os commits
incluídos nas releases. A paginação da API evita que essas comparações sejam
limitadas silenciosamente a uma fração dos commits \cite{github}.

A validade dos proxies será examinada por três avaliadores independentes em
60 repositórios, com cinco releases por projeto. A concordância e a qualidade
da heurística corretiva fornecerão evidências para interpretar as métricas.
A comparação entre combinações de releases, pré-releases, tags e variantes
de lead time e CFR avaliará a estabilidade das classificações.

\subsection{Hipóteses informais propostas}

As hipóteses abaixo constituem uma proposta anterior à coleta deste estudo e
devem ser aprovadas pelo trio no protocolo. Elas expressam expectativas a
confrontar com os dados, não resultados observados.

\begin{description}
\item[H01 (RQ01).] Espera-se heterogeneidade na frequência de releases,
com concentração abaixo de uma entrega semanal, pois projetos podem
agrupar mudanças em ciclos distintos de publicação.
\item[H02 (RQ02).] Espera-se que a variante por release apresente, em geral,
lead times superiores aos da variante por commit, pela influência de commits
antigos. Não é uma identidade matemática: as variantes ponderam as releases
de formas diferentes.
\item[H03 (RQ03).] Espera-se divergência entre CFR baseado em CI e baseado
em releases corretivas, pois capturam eventos distintos. Não se antecipa
ordenação obrigatória entre seus valores.
\item[H04 (RQ04).] Espera-se distribuição assimétrica dos tempos de
recuperação, com muitos episódios curtos e alguns prolongados. Episódios
não encerrados durante a janela deverão permanecer registrados como censurados.
\item[H05 (RQ05).] Espera-se associação monotônica fraca ou não positiva
entre frequência e CFR. Spearman será calculado para ambas as variantes
de CFR, sem interpretar associação como causalidade.
\item[H06 (RQ06).] Espera-se diferença entre subgrupos em pelo menos uma
métrica. Propõem-se linguagem principal, quartis de estrelas e quartis de
contribuidores como fatores, sujeitos à aprovação antes de observar resultados.
\item[H07 (RQ07).] Espera-se que parte dos repositórios mude de categoria
DORA ao alterar as definições de entrega, lead time e CFR, especialmente
próximo aos cortes de classificação.
\end{description}

Os artefatos serão organizados em módulos de coleta, normalização, métricas e
análise, com testes e rastreabilidade por issues. Validação manual, sensibilidade
e replicação cruzada deverão sustentar a discussão dos limites do estudo.
Repositório: \url{https://github.com/lucasrsnd/lab-medicao-experimentacao}.
Acompanhamento: \url{https://github.com/users/lucasrsnd/projects/2}.

